Flooding is among the costliest natural hazards in the conterminous United States (CONUS)~\cite{fema_cost_of_flooding}, with damage dynamics that vary sharply by topography, drainage infrastructure, and climate regime~\cite{brunner_spatial_2020}.
%
Among flood types, \emph{pluvial} flooding (inundation driven by intense rainfall exceeding local drainage capacity) is a particularly difficult prediction target~\cite{rosenzweig_pluvial_2018,nelson-mercer_pluvial_2025}.
%
% Events are highly localized in space and time, develop on hourly timescales, and account for roughly $45\%$ of National Flood Insurance Program (NFIP) claims and approximately \$56~billion of insured damage since 1980~\cite{fema_nfip_claims_v2}.
%
Anticipating pluvial damage at a resolution operationally meaningful to insurers, emergency managers, and municipal planners is both an open scientific problem and is  of immediate societal value.

Despite this urgency, pluvial flood damage prediction lags behind its riverine and coastal counterparts~\cite{rosenzweig_pluvial_2018}.
%
Riverine flooding benefits from a dense network of stream gauges that provide a relatively clean, continuous training signal~\cite{boutayeb_when_2025,nearing_global_2024}, and coastal flooding is similarly well-tracked by tide and surge gauges~\cite{bates_combined_2021,yang_predicting_2022}.
%
Pluvial flooding, on the other hand, has a noisy observational backbone, and its damage signal lives in scattered insurance claims~\cite{fema_nfip_claims_v2} and citizen reports~\cite{puttinaovarat_flood_2020,mayo_groundsource_2026}, each carrying nontrivial spatial and temporal uncertainty.
%
CONUS-scale pluvial products do exist, but they occupy a different operating regime than ours.
%
FEMA flood maps~\cite{fema_flood_zone_2021} are static delineations of long-term flood hazard zones rather than dynamic estimates of damage, while CONUS-scale hydrodynamic models~\cite{wing_30_2024,bates_combined_2021} simulate physical inundation at high fidelity but are computationally heavy and produced for fixed design scenarios rather than on a daily basis.
%
Neither delivers a \emph{lightweight, daily, ${\sim}1$\,km prediction of insured pluvial damage}, the regime we target.

We close this gap with \textbf{DELUGE} (\textit{\textbf{D}}aily \textit{\textbf{E}}stimation of p\textit{\textbf{LU}}vial flood dama\textit{\textbf{GE}}), to our knowledge the first end-to-end learned system to predict daily, ${\sim}1$\,km insured pluvial flood damage across the highest-claim regions of CONUS through interpretable conditioning on Earth-observation foundation-model embeddings.
%
DELUGE is a multimodal CNN whose local inductive bias matches the locally driven nature of pluvial dynamics, in contrast to the catchment-scale connectivity behind GNN-based riverine approaches~\cite{sarkar_hydrogat_2025,kazadi_floodgnn-gru_2024}.
%
We train DELUGE on NFIP claims as a damage proxy, after applying temporal and spatial uncertainty corrections that refine each claim's timing and footprint.
%
Under a spatial block holdout protocol that partitions 75\,km grid cells into a train/test split, DELUGE achieves a PR-AUC of \tsim$0.24$, well above the \tsim$0.0025$ no-skill baseline implied by the task's \tsim$0.25\%$ positive rate, and outperforms tuned Random Forest, XGBoost, and LightGBM baselines commonly used in claims-based flood-damage prediction, with the largest gains on high-cost claims.
% (standard in the pluvial flood literature) by 9\% to 26\% on PR-AUC and 11\% to 35\% on a dollar-weighted variant.

Beyond its predictive performance, DELUGE's architectural novelty lies in two parametric modules within the hydrometeorology branch.
%
The \emph{Value Modulator} learns a per-sample monotone warp of each hydrometeorological input, and the \emph{Temporal Modulator} learns a per-sample temporal-response kernel that causally convolves the hydrometeorological time series.
%
Both modules are conditioned on local place characteristics, encoded through terrain descriptors and Google AlphaEarth foundation-model embeddings~\cite{brown_alphaearth_2025}, which supply a dense, learned representation of the built and natural environment that varies continuously across CONUS.
% and captures local context that hand-crafted terrain descriptors alone miss.

Rather than concatenating these embeddings into a black-box encoder, we route them through the physics-shaped modulators, whose parameters carry direct hydrological meaning such as response peak lag in hours and decay timescale.
%
This makes the modulators a principled and interpretable way to consume an otherwise opaque foundation embedding, yielding \emph{interpretability-by-design}~\cite{rudin_stop_2019} where ConvLSTM or Transformer-style encoders would require post-hoc attribution approximations like SHAP~\cite{lundberg_unified_2017}.
%
Thus our interpretability is at the architecture level, exposing the per-location hydrological response DELUGE has learned, rather than a per-prediction explanation.
%
Our main contributions are:
\begin{enumerate}[nosep, leftmargin=14pt]
    \item \textbf{DELUGE}, the first end-to-end learned system for daily pluvial flood damage prediction at ${\sim}1$\,km across the highest-claim regions of CONUS, outperforming tree-based baselines standard in the pluvial flood prediction literature, with the largest gains on high-cost claims.
    \item Two uncertainty-correction procedures for NFIP claims, a precipitation guided correction of claim dates and a spatial refinement that intersects redacted geographies to localize damage to higher resolution.
    \item The \emph{Value Modulator} and \emph{Temporal Modulator}, terrain- and AlphaEarth-conditioned parametric modules whose learned per-patch parameters are themselves hydrologically meaningful, yielding \emph{interpretability-by-design} of Geospatial Foundation Model integration.
    \item Inspection of the learned modulator parameters via geography-blind clustering, showing they recover hydrologically coherent regimes consistent with physical expectation.
    % \item Empirical verification of the modulators' physical fidelity via geography-blind clustering that recovers physically grounded explanation of identified regimes.
\end{enumerate}
\vspace{-5pt}

% DELUGE shows that with appropriate parametric inductive biases, sparse and noisy flood damage labels are enough to recover physically meaningful structure, supporting future data-driven approaches to continental-scale flood impact modeling.
